%% ma: L12-B08-0173
%% lop: 12
%% chuong: 2
%% bai: 8
%% dang: 12.8.53
%% loai: TN4
%% muc: VD
%% dap_an: "A"
%% nguon: "CĐ2.4.pdf giáo viên gửi 27/10/2026 (M07, Lê Minh Khôi – Nguyễn Anh Khôi, Chương II Chủ đề 4 Các bài toán về tâm tỉ cự), Câu 23"
%% kien_thuc: "Bài 8 (khoảng cách hai điểm, biểu thức tọa độ), Bài 6 (bình phương vô hướng, tâm tỉ cự)"
%% trang_thai: chinh-thuc
%% ghi_chu: "Giáo viên đã duyệt 27/10/2026"
\begin{ex}
Trong không gian $Oxyz$, cho ba điểm $A(2;0;3)$, $B(-1;3;0)$, $C(1;1;2)$. Với $M$ là điểm bất kỳ trong không gian, tìm giá trị lớn nhất của biểu thức $S=2MA^2+MB^2-12MC^2$.
\choice{\True $S_{\max}=18$}{$S_{\max}=4$}{$S_{\max}=10$}{$S_{\max}=15$}
\loigiai{Gọi $I$ thỏa $2\overrightarrow{IA}+\overrightarrow{IB}-12\overrightarrow{IC}=\vec 0\Rightarrow I=\dfrac{2A+B-12C}{-9}=(1;1;2)$.

$S=-9MI^2+2IA^2+IB^2-12IC^2\le2\cdot3+12-0=18$.}
\end{ex}
